\section{Critérios de sucesso}
\label{sec:criterios_sucesso}

Os critérios de sucesso estabelecem as condições mínimas que definem o êxito da operação de RVD/B.

\subsection{Primários}

\begin{enumerate}[label=(\roman*), leftmargin=*, itemsep=0pt]
\item \textbf{Conclusão do acoplamento:} o visitante deve concluir o acoplamento dentro da janela temporal definida para a operação.

\item \textbf{Sem violação de segurança:} a trajetória da aproximação curta deve respeitar integralmente as zonas de segurança estabelecidas, conforme:
    \begin{enumerate}[label=(\alph*), itemsep=0pt]
    \item nenhuma entrada na \textit{keep-out zone} (KOZ);  
    \item permanência contínua dentro do corredor de aproximação durante toda a aproximação curta.
    \end{enumerate}

\item \textbf{Taxa de aproximação no contato:} a velocidade relativa entre o visitante e o alvo no momento de contato deve permanecer abaixo do limite máximo estabelecido para o acoplamento.

\item \textbf{Alinhamento no acoplamento:} o visitante deve atingir o contato mantendo a orientação e o posicionamento dentro dos limites de alinhamento especificados, conforme:
    \begin{enumerate}[label=(\alph*), itemsep=0pt]
    \item erro relativo de orientação dentro da margem aceitável;  
    \item desalinhamento lateral inferior aos limites definidos para a interface mecânica.
    \end{enumerate}
\end{enumerate}

\subsection{Desempenho de guiagem e controle}

Avaliou-se os erros relativos de posição, velocidade e atitude e a coerência das comutações entre modos.

\begin{enumerate}[label=(\roman*), leftmargin=*, itemsep=0pt]
\item \textbf{Erro de posição e velocidade relativa:} analisado por fase, considerando as métricas de valor médio, quadrático médio (RMS) e valor máximo registrado durante cada regime de aproximação.

\item \textbf{Erro de atitude relativa e taxa angular:} a atitude do visitante em relação ao alvo deve manter-se estável durante a aproximação curta e final, com variações angulares dentro das margens pré-determinadas para o alinhamento.

\item \textbf{Número de comutações de modo:} as transições entre os diferentes ganhos do controlador PD e modos de guiagem devem ocorrer conforme a lógica definida, sem oscilações indesejadas ou comutação excessiva.
\end{enumerate}

\subsection{Recursos}

Os critérios aqui apresentados avaliam a eficiência do sistema quanto ao consumo de propelente e ao tempo de execução das manobras.

\begin{enumerate}[label=(\roman*), leftmargin=*, itemsep=0pt]
\item \textbf{$\Delta v$ total:} é a soma de todos os incrementos de velocidade aplicados durante as fases de transferência e aproximação.

\item \textbf{Margem de propelente remanescente:} representa a diferença entre o propelente disponível após o acoplamento e a quantidade inicial de propelente.
\end{enumerate}

\subsection{Conformidade operacional e segurança}

Os critérios de conformidade operacional e segurança avaliam se a missão foi conduzida dentro dos limites definidos para cada fase da aproximação.

\begin{enumerate}[label=(\roman*), leftmargin=*, itemsep=0pt]
\item \textbf{Violações de limites por fase:} devem ser analisadas separadamente para:
    \begin{enumerate}[label=(\alph*), itemsep=0pt]
    \item velocidade relativa;  
    \item atitude;  
    \item saída do corredor de aproximação.
    \end{enumerate}

\item \textbf{Registro de violações:} cada evento deve ser documentado quanto à quantidade e à duração, sendo que ambos os indicadores devem ser nulos para que a manobra seja considerada aprovada.
\end{enumerate}

\subsection{Manipulador robótico}

Devido à arquitetura 4R1P adotada, o desempenho do manipulador robótico é avaliado em respeito ao acoplamento, para que seja realizado de forma precisa, sem esforços excessivos sobre a interface mecânica.

\begin{enumerate}[label=(\roman*), leftmargin=*, itemsep=0pt]
\item \textbf{Posicionamento da ferramenta:} o erro de pose da ferramenta na fase de pré-captura deve permanecer dentro da margem de tolerância definida, cuidando o alinhamento entre a junta prismática e a porta de acoplamento.

\item \textbf{Força e torque de contato:} durante o engate, as forças e torques aplicados pela ferramenta devem permanecer dentro dos limites operacionais especificados para a interface.

\item \textbf{Tempo de estabilização pós-contato:} logo após o engate, o acoplamento entre base--manipulador deve ser estável dentro do tempo previsto e sem oscilações perceptíveis.
\end{enumerate}
